% Source-derived chapter 07: The divisorial contraction case
% Original source: riemann-roch-polynomials-hyperkahler-fourfolds
\chapter{The divisorial contraction case}
\label{riemann-roch-polynomials-hyperkahler-fourfolds-chapter-07}
\source{riemann-roch-polynomials-hyperkahler-fourfolds}

\begin{remark}[Source section]
\label{riemann-roch-polynomials-hyperkahler-fourfolds-chapter-07-source}
\source{riemann-roch-polynomials-hyperkahler-fourfolds}
\source{riemann-roch-polynomials-hyperkahler-fourfolds}
This chapter reproduces the corresponding section of the retrieved
LaTeX source, including its definitions, statements, examples, and
proofs. It has not yet been translated into Lean.
\notready
\end{remark}

% The source section title was: The divisorial contraction case
\label{sec:DivisorialContraction}
\source{riemann-roch-polynomials-hyperkahler-fourfolds}

Let $(X,\lll,\mm)$ be a very general triple satisfying~\eqref{hypo}.\
By Proposition~\ref{prop:SurfacesOnHK4}, we have  $\NS(X)=\Z\lll\oplus\Z\mm$.\
We assume in this section that we are in case~\ref{caseC2}: the class $\lll$ is nef, while the class~$\mm$ is isotropic and not nef.\ Thus, we have  $\Mov(X)=\R_{\ge 0}\lll+\R_{\ge 0}(\lll+\mm)$ and there is a divisorial contraction  $c\colon X\to Y$ defined by  some positive power of the semi-ample line bundle $L\otimes M$, whose exceptional locus is the prime divisor $E$ with class $-\lll+\mm$.\
Our aim is to prove Proposition~\ref{prop:CaseC2}: some tensor power of $L$ defines a Lagrangian fibration on $X$.

The fourfold $Y$ is Gorenstein with trivial canonical sheaf, rational singularities, and singular locus the surface $\Sigma\coloneqq c(E)$.\
The class in $H_2(X,\Z)\isom H^2(X ,\Z)^\vee$ of a general fiber of $E\to \Sigma$ is given by the linear form $q_X(-\lll+\mm,\bullet)$ which, since   $q_X(-\lll+\mm,\lll)=1$, is nondivisible.\
The class of any curve contracted by $c$ is a multiple of that class, hence all $1$-dimensional fibers of $c$ are irreducible, smooth rational curves.

\begin{prop}
\source{riemann-roch-polynomials-hyperkahler-fourfolds}
\notready\label{prop:NoFibersDimension2}
\source{riemann-roch-polynomials-hyperkahler-fourfolds}
The fibers of the contraction $c\colon X\to Y$ all have dimension at most~$ 1$, the varieties $E$ and $\Sigma$ are smooth, and the restriction $c_E:=c\vert_E\colon E\to \Sigma$ is a $\P^1$-fibration.
\end{prop}

\begin{proof}
We argue by contradiction and assume that there is an integral surface $S\subset X$ such that $c(S)$ is a point in $Y$.\ Then $S\subset E$.\ We first claim that the class
\begin{equation}\label{eq:S'}
\source{riemann-roch-polynomials-hyperkahler-fourfolds}
[S'] \coloneqq 2(\lll+\mm)(-\lll+\mm) - [S]
\end{equation}
in ${\Hdg}^4(X,\Z)$ is effective.

Indeed, by assumption, the line bundles $L^2\otimes M^2$ and $L^3\otimes M$ are big and nef on $X$.\ By the Kawamata--Viehweg vanishing theorem and~\eqref{hypoprime}, we get
\[
h^0(X, L^2\otimes M^2)=21\quad\text{ and }\quad h^0(X, L^3\otimes M)=15.
\]
It follows that the restriction map
\[
H^0(X,L^2\otimes M^2)\rightarrow H^0(E,(L^2\otimes M^2)\vert_E),
\]
whose kernel equals $H^0(X,L^3\otimes M)$, has rank at least~$6$.\
On the other hand, as $L\otimes M$ is numerically trivial on $S$, the restriction map
\[
H^0(X,L^2\otimes M^2)\rightarrow H^0(S,(L^2\otimes M^2)\vert_S)
\]
has rank at most~$1$.\
It follows that  there exists a section of $L^2\otimes M^2$ vanishing on $S$ which does not vanish identically on $E$.\
Since the class of $E$ is $-\lll+\mm$, the claim follows.

We now compute the  intersection matrices $M_{[S]}$ and $M_{[S']}$ defined in~\eqref{eqmatriceO}.\ They are nonzero since there is an ample class  in $\Z\lll\op\Z\mm$.\
Since the surface $S$ is contracted by $c$, the line bundle $L\otimes M$ is numerically trivial on $S$ and we get
\[
M_{[S]}=
\begin{pmatrix}
t & -t \\
-t & t
\end{pmatrix}
\]
for some positive integer $t$.\
Hence, using~\eqref{eq:S'}, we get
\[
M_{[S']}=
\begin{pmatrix}
4-t & t \\
t & -4-t
\end{pmatrix}.
\]
Since $[S']$ is effective, we get $4-t\geq0$, hence $t\in\{1,2,3,4\}$.

By Proposition~\ref{prop:SurfacesOnHK4}, we can write
\begin{equation}\label{eq:Paris20211103}
\source{riemann-roch-polynomials-hyperkahler-fourfolds}
\begin{split}
    &[S] = \frac t2\, \Big( \lll^2 -\lll\mm + \mm^2 \Big) +   w\,\Big(q_X^\vee - \frac{25}{2} \lll\mm \Big), \\
    &[S'] = -\Big(2+ \frac t2 \Big)\lll^2 + \frac t2 \lll\mm + \Big(2 - \frac t2 \Big)\mm^2 -   w\,\Big(q_X^\vee - \frac{25}{2} \lll\mm \Big),
\end{split}
\end{equation}
for some $w\in\Q$.\
Since $[S]$ is an integral class, we have $[S]^2\in\Z$.\
Using~\eqref{qxc}, we get
\[
2[S]^2 =   t^2 + 525\, w^2=t^2 + 3\cdot  7\, (5w)^2,
\]
so that $5w\in\Z$.

From~\eqref{eq:Paris20211103}, by the same reasoning used at the end of the proof of Lemma~\ref{lecasholhoM}, for any class $\omega\in \overline{\cK}_X\subset H^{1,1}(X,\R)$ (in our situation this means in particular that $q_X(\omega,\lll)\ge0$ and $q_X(\omega,-\lll+\mm)\geq0$) with $q_X(\omega)=0$, we obtain
\begin{equation}\label{eq:Paris20211103_second}
\source{riemann-roch-polynomials-hyperkahler-fourfolds}
\begin{split}
    & 0\leq \int_X [S]\omega^2 = \frac t2 \int_X\lll^2\omega^2- \Big(\frac {t+25w}2\Big)\int_X \lll\mm\omega^2 + \frac t2\int_X \mm^2\omega^2,  \\
    & 0\leq \int_X [S']\omega^2 = - \Big(2+\frac t2\Big)\int_X \lll^2\omega^2
     + \Big(\frac {t+25w}2\Big)\int_X \lll\mm\omega^2
     +\Big(2-\frac t2\Big)\int_X \mm^2\omega^2 .
\end{split}
\end{equation}
Moreover, since $q_X(\omega)=0$ and $c_X=3$,   the Fujiki relation~\eqref{eq:Fujiki3} implies, for all $\alpha,\beta\in H^2(X,\Z)$,
\[
\int_X \alpha\beta\omega^2 = 2 q_X(\alpha,\omega)\, q_X(\beta,\omega).
\]
Thus, from~\eqref{eq:Paris20211103_second}, we deduce
\[
\frac {t+25w}2  \leq t\quad\text{ and }\quad \frac {t+25w}2  \geq t,
\]
and thus
\[
25 w = t, \qquad \text{ with } t\in\{1,2,3,4\}\text{ and }5w\in\Z,
\]
which is impossible.\ This proves that $c$ contracts no surfaces.

The other statements of the lemma then follow from
\cite[Theorem~1.3(ii)]{wie}.
\end{proof}

Since the line bundle $L$ has intersection~1 with all fibers of $c_E$, we immediately deduce the following result.

\begin{lemm}\label{lem:EisP1bundle}
\source{riemann-roch-polynomials-hyperkahler-fourfolds}
Let $\cE\coloneqq c_{E\,*}L$.
We have an isomorphism $E\cong\P_{\Sigma}(\cE^\vee)$ over $\Sigma$.
\end{lemm}

By Lemma~\ref{lem:EisP1bundle}, the line bundle $(L\otimes M)\vert_E$ descends to a line bundle $H_{\Sigma}$ on~$\Sigma$, which is ample since some positive power of $L\otimes M$ is the pullback of an ample line bundle on $Y$.\
We will study  in more detail in Section~\ref{sec:OGrady} the surface $\Sigma$, the polarization $H_{\Sigma}$, and the rank-2 vector bundle $\cE$.
For the moment, we just use their existence to finish the proof of Proposition~\ref{prop:CaseC2}.
We need one last result before that.

\begin{lemm}\label{levanh2}
\source{riemann-roch-polynomials-hyperkahler-fourfolds}
One has $H^2(X, M^{\otp 2})=H^4(X, M^{\otp 2})=0$ and $h^0(X, M^{\otp 2})\ge3$.
\end{lemm}

\begin{proof}
We first prove $H^2(X, M^{-2})=0$.\
Consider the exact sequence
\begin{equation*}
0\to \cO_X(-E)\to \cO_X\to \cO_E\to 0.
\end{equation*}
Tensoring by $(L\otimes M)^{\otp -1}$, we obtain
\begin{equation*}
0\to  M^{-2}\to (L\otimes M)^{\otp -1}\to (L\otimes M)^{\otp -1}\vert_E\to 0.
\end{equation*}
Since $H^2(X,(L\otimes M)^{\otp -1})=H^2(X,L\otimes M)=0$, it  suffices to prove  $H^1(E,(L\otimes M)^{\otp -1}\vert_E)=0$.\
But, by Lemma~\ref{lem:EisP1bundle}, we have
\[
H^1(E,(L\otimes M)^{-1}\vert_E)\isom H^1(\Sigma, H_{\Sigma}^{-1})=0,
\]
since $H_{\Sigma}$ is ample on the smooth surface $\Sigma$, as we wanted.

By Serre duality, we get $H^2(X,M^{\otp 2})=0$.\
We also have $H^4(X,M^{\otp 2})=H^0(X,M^{\otp -2})=0$.

Finally, the Riemann--Roch theorem takes the form
\begin{equation}\label{eq:Paris20211103_third}
\source{riemann-roch-polynomials-hyperkahler-fourfolds}
h^0(X,M^{\otp  2})+h^2(X,M^{\otp  2})+h^4(X,M^{\otp2})\geq \chi(X,M^{\otp  2})=P_{RR,X}(0)=3,
\end{equation}
completing the proof of the lemma.
\end{proof}

In fact, the divisor $E$ is fixed in $|M^{\otp 2}|$, hence $h^0(X,M^{\otp  2})= h^0(X,L\otimes M)=6$, but we will use again the Riemann--Roch argument~\eqref{eq:Paris20211103_third} in the proof below.

\begin{proof}[Proof of Proposition~\ref{prop:CaseC2}.]
As observed in Section~\ref{subsec:Lagrangian4folds}, we only have to show  $h^0(X,L^{\otp 2})\geq 2$.\
This inequality follows from Lemma \ref{levanh2} by deformation and specialization: as we explained in Section~\ref{subsec:ExceptionalDivisor}, the reflection that permutes $\lll$ and $\mm$ is a monodromy operator.\
This means that one can deform the pair $(X,L)$ into the pair $(X,M)$ through a family $(\cX,\cL)\to T$.\
By semi-continuity, Lemma~\ref{levanh2} implies $h^2(\cX_t, \cL_t^{\otp 2})=h^4(\cX_t, \cL_t^{\otp 2})=0$ for $t\in T$ general.\
The Riemann--Roch argument of~\eqref{eq:Paris20211103_third} then implies $h^0(\cX_t, \cL_t^{\otp 2})\geq 3$ for $t\in T$ general.\
By upper semi-continuity again, we conclude $h^0(X,L^{\otp 2})\geq 3$, as we wanted.
\end{proof}


%%%%%%%%%%%%%%%%%%%%%
